\documentclass[11pt]{article}

\usepackage[margin=1in]{geometry}
\usepackage{amsmath,amssymb,amsthm}
\usepackage{graphicx}
\usepackage{hyperref}
\hypersetup{hidelinks}

\newtheorem{theorem}{Theorem}
\newtheorem{lemma}[theorem]{Lemma}
\newcommand{\A}{\mathcal A}
\newcommand{\Mcb}{M_{\mathrm{cb}}}
\newcommand{\otA}{\widehat\otimes_{\A}}

\title{Completely bounded multipliers of a Beurling--Fourier algebra\\
form a dual Banach algebra}
\author{Run \texttt{fa\_banach\_001}, lane 19}
\date{August 11, 2026}

\begin{document}
\maketitle

\begin{abstract}
Alaghmandan, Giselsson, Samei, and Turowska prove that, for every locally
compact group $G$ and weight inverse $\omega$, the completely bounded
multiplier algebra $\Mcb A(G,\omega)$ is completely isometric to the dual of
the balanced operator-projective tensor product
$A(G,\omega)\widehat\otimes_{A(G,\omega)}C_r^*(G)$.  They state that they
believe this multiplier algebra is a dual Banach algebra, proving that fact
only when $G$ is amenable.  We prove the general statement.  For a fixed
multiplier $\psi$, its multiplication map on $A(G,\omega)$ preserves the
balanced tensor relations and therefore induces a completely bounded map
$S_\psi$ on the displayed predual.  The adjoint $S_\psi^*$ is exactly
multiplication by $\psi$ on $\Mcb A(G,\omega)$.  Multiplication is consequently
separately weak-star continuous for the source predual.
\end{abstract}

\section{Source statement}

The source is M. Alaghmandan, O. Giselsson, E. Samei, and L. Turowska,
\emph{Multipliers of Beurling--Fourier algebras}, arXiv:2606.06675v1
\cite{Source}.  On page 27, immediately before Theorem 7.12, the authors write:
\begin{quote}
``It is known that $\Mcb A(G)$ is a dual Banach algebra.  We believe that
this should hold for $\Mcb A(G,\omega)$ as well.  By Theorem 8.3 and Lemma
5.5, this is true for amenable groups.''
\end{quote}
The full page-width source passage, including the balanced tensor definition,
the conjectural sentence, and Theorem 7.12, is reproduced below.

\begin{center}
\includegraphics[width=0.98\textwidth]{figures/open_problem_crop.png}
\end{center}

Theorem 7.12 of the source proves, for arbitrary locally compact $G$ and every
weight inverse $\omega$, the completely isometric identification
\[
 \Mcb A(G,\omega)
 \cong
 \left(A(G,\omega)\widehat\otimes_{A(G,\omega)}C_r^*(G)\right)^*.
\]
The unresolved point is whether multiplication is separately weak-star
continuous for this predual.

\section{Proof intuition}

The predual constructed by the source already remembers the multiplier
module structure.  Indeed, let $\psi$ be a fixed completely bounded
multiplier.  Applying $\psi$ to the first tensor factor sends
\[
 ab\otimes x-a\otimes b\cdot x
 \quad\hbox{to}\quad
 \psi(ab)\otimes x-(\psi a)\otimes b\cdot x.
\]
The multiplier identity $\psi(ab)=(\psi a)b$ says that the image is another
balanced relation.  Thus the map descends to the quotient predual.  Under the
source pairing, its adjoint sends a multiplier $\phi$ to the product
$\phi\psi$.  This gives the missing preadjoint of multiplication directly;
no approximation, amenability, or additional property of the weight is
needed.

\section{An abstract tensor lemma}

We use the operator-space projective tensor product and its standard
functoriality \cite{BLM}.  The proof below also records explicitly the only
quotient issue that matters.

\begin{lemma}\label{lem:abstract}
Let $\A$ be a commutative completely contractive Banach algebra, let $X$ be a
left operator $\A$-module, and put
\[
 P=\A\widehat\otimes_{\A}X.
\]
Suppose $\Mcb(\A)$ is identified with $P^*$ through the canonical pairing
\[
 \langle\phi,[a\otimes x]\rangle=\langle\phi(a),x\rangle.
\]
Then this identification makes $\Mcb(\A)$ a dual Banach algebra.  More
precisely, for every $\psi\in\Mcb(\A)$, multiplication by $\psi$ is the
adjoint of a completely bounded operator on $P$.
\end{lemma}

\begin{proof}
Write $R_\psi:\A\to\A$ for the completely bounded multiplier map
$R_\psi(a)=\psi(a)$.  By functoriality,
$R_\psi\widehat\otimes\mathrm{id}_X$ is completely bounded on
$\A\widehat\otimes X$.

Let $N$ be the closed balanced subspace generated by
\[
 ab\otimes x-a\otimes b\cdot x
 \qquad(a,b\in\A,\ x\in X).
\]
Since $R_\psi$ is a right $\A$-module map,
\[
 R_\psi(ab)=R_\psi(a)b.
\]
Consequently
\begin{align*}
 (R_\psi\widehat\otimes\mathrm{id}_X)
 (ab\otimes x-a\otimes b\cdot x)
 &=R_\psi(ab)\otimes x-R_\psi(a)\otimes b\cdot x\\
 &=(R_\psi(a)b)\otimes x-R_\psi(a)\otimes b\cdot x
 \in N.
\end{align*}
Thus $R_\psi\widehat\otimes\mathrm{id}_X$ descends to a completely bounded
operator $S_\psi:P\to P$ satisfying
\[
 S_\psi[a\otimes x]=[\psi(a)\otimes x].
\]

For $\phi\in\Mcb(\A)$ and an elementary balanced tensor,
\begin{align*}
 \langle S_\psi^*\phi,[a\otimes x]\rangle
 &=\langle\phi,S_\psi[a\otimes x]\rangle\\
 &=\langle\phi(\psi(a)),x\rangle\\
 &=\langle(\phi\psi)(a),x\rangle\\
 &=\langle\phi\psi,[a\otimes x]\rangle.
\end{align*}
Elementary tensors are dense in $P$, so $S_\psi^*\phi=\phi\psi$.  Hence
multiplication by each fixed $\psi$ is weak-star continuous.  Since $\A$ is
commutative, this proves separate weak-star continuity of multiplication.
\end{proof}

\section{Resolution of the source conjecture}

\begin{theorem}\label{thm:main}
Let $G$ be a locally compact group and let $\omega$ be a weight inverse on
$G$.  With the predual from Theorem 7.12 of the source,
\[
 \Mcb A(G,\omega)
 =\left(A(G,\omega)
 \widehat\otimes_{A(G,\omega)}C_r^*(G)\right)^*
\]
is a dual Banach algebra.  In fact, multiplication is separately completely
weak-star continuous.
\end{theorem}

\begin{proof}
Set $\A=A(G,\omega)$ and $X=C_r^*(G)$, with the left $\A$-module action
used in the source balanced tensor product.  The algebra $\A$ is a
commutative completely contractive Banach algebra, and a member of
$\Mcb A(G,\omega)$ is precisely a completely bounded multiplier of $\A$.
Theorem 7.12 of the source supplies the completely isometric dual
identification and the canonical pairing required in
Lemma~\ref{lem:abstract}.  Applying that lemma proves the assertion.  The
maps $S_\psi$ are completely bounded, so their adjoints are completely
weak-star continuous at every matrix level.
\end{proof}

\section{Verification and novelty notes}

The argument depends only on the source's Theorem 7.12 and on functoriality of
the operator-projective tensor product.  The balanced quotient is checked in
line.  The key identity can be audited on one elementary tensor:
\[
 \langle\phi\psi,[a\otimes x]\rangle
 =\langle\phi,[\psi(a)\otimes x]\rangle.
\]
This identifies multiplication with an adjoint operator before taking any
limit, so no density-of-products argument or uniform-boundedness step is
needed.

A bounded novelty check on August 11, 2026 searched all four lightweight run
indexes for arXiv:2606.06675, Beurling--Fourier multipliers, and dual Banach
algebras.  External searches covered the exact title and arXiv id and the
phrases ``$M_{\rm cb}A(G,\omega)$ dual Banach algebra'', ``Beurling--Fourier
dual Banach algebra'', and ``balanced tensor multiplier predual''.  The arXiv
record still has only v1, submitted June 4, 2026.  No separate answer or
matching general theorem was found.  Novelty remains subject to specialist
review.

\section{Limitations and human-review recommendation}

The result uses exactly the weak-star topology induced by the source predual;
it does not assert uniqueness of that predual.  It assumes Theorem 7.12 of the
source as stated.  Within that framework, there is no remaining conditional
step and no restriction to amenable groups.

This packet is recommended for expert review as a likely valid full solution.
The most important checks are: (i) that the first tensor factor in the source
balanced product is the right regular $A(G,\omega)$-module, (ii) that a
completely bounded multiplier is a completely bounded right module map, and
(iii) that the source pairing agrees with the pairing used in
Lemma~\ref{lem:abstract}.

\begin{thebibliography}{9}
\bibitem{Source}
M. Alaghmandan, O. Giselsson, E. Samei, and L. Turowska,
\emph{Multipliers of Beurling--Fourier algebras},
arXiv:2606.06675v1, 2026.

\bibitem{BLM}
D. P. Blecher and C. Le Merdy,
\emph{Operator Algebras and Their Modules: An Operator Space Approach},
Oxford University Press, 2004.
\end{thebibliography}

\end{document}
